\section{Data Structures}

\topic{Lazy Segment Tree}[\log n]

{\small\color{muted} Life will be a lot easier if you read this after understanding the basic segment tree.\par}

\begin{problem}
Given an array \(a\) of length \(n\), we can do the following in \(\bigO{\log n}\).
\begin{itemize}
\item Update a whole range \([l, r)\). By update we mean an operation done to every element of the range, where the operation, written \(x \circ f\) for the update \(f\) applied to \(x\), must satisfy the following.
\begin{enumerate}[label=\textup{(U\arabic*)},leftmargin=*]
\item It can be applied to a whole range at once.
\item Updates are closed under composition, \((x \circ g) \circ f = x \circ h\) for an update \(h\).
\item Updates have an identity \(\mathrm{id}\), with \(x \circ \mathrm{id} = x\).
\end{enumerate}
\item Query a range \([l, r)\). By query we mean an operation that combines all elements of the range, where the operation, written \(x \oplus y\), must satisfy the following.
\begin{enumerate}[label=\textup{(Q\arabic*)},leftmargin=*]
\item It is associative, \((x \oplus y) \oplus z = x \oplus (y \oplus z)\).
\item It has an identity \(e\), with \(e \oplus x = x \oplus e = x\).
\end{enumerate}
\end{itemize}
\end{problem}

\begin{notes}[Reading the definition]
Remember that in a segment tree, each node holds the value of one range.
\begin{enumerate}[label=\textup{(U\arabic*)},leftmargin=*]
\item When we arrive at the node of a range \([l, r)\) that we want to update, we can update it using only its old value, the update and often the length of the range. Say the query is sum and we want to add \(c\) to \([l, r)\). \\
After walking down the tree, we arrive at the node of \([l, r)\). \\
We do not go deeper to update every leaf. \\
We simply add \(c \cdot (r - l)\) to its value, as if \(c\) were added to every element below. \\
Other updates with this property are the following.
\begin{itemize}
\item Setting every element to \(c\).
\item Multiplying by \(c\) when the query is sum.
\item Adding \(c\) when the query is minimum.
\end{itemize}
It is a good exercise to figure out how to apply each one.
\item A node may get a second update before the first one reaches its children, and then the two must become one. Some updates that satisfy this property are the following.
\begin{itemize}
\item Adding. Adding \(c\) and then \(d\) is the same as adding \(c + d\).
\item Multiplying. Multiplying by \(c\) and then by \(d\) is the same as multiplying by \(cd\).
\item Setting. Setting to \(c\) and then to \(d\) is the same as setting to \(d\).
\item Multiplying and adding. \\
Doing \(x \mapsto ax + b\) and then \(x \mapsto cx + d\) is the same as \(x \mapsto acx + bc + d\). \\
This is like \(c(ax + b) + d\). So any operations of this type can be merged into one.
\end{itemize}
Other updates with this property are the following.
\begin{itemize}
\item XOR with \(c\).
\item Setting and adding mixed together.
\item Taking the maximum with \(c\).
\end{itemize}
It is a good exercise to figure out how to merge each one.
\end{enumerate}
\end{notes}

\subsubsection{Tree Structure}
We draw the segment tree as a stack of bricks. Every node is a brick exactly as wide as its range, and it sits on top of its two halves. The root covers \([0, n)\). A node covering \([\mathit{lo}, \mathit{hi})\) splits at \(m = \floor{(\mathit{lo} + \mathit{hi}) / 2}\) into \([\mathit{lo}, m)\) and \([m, \mathit{hi})\), and its value is the query of its range. The bricks of width one are the leaves, one per element. From here on the query is sum.

\begin{center}
\begin{tikzpicture}[x=6.5mm,y=6.5mm]
  \brick{0}{6}{3}{23}
  \brick{0}{3}{2}{8}
  \brick{3}{6}{2}{15}
  \foreach \lo/\hi/\v in {0/1/3, 1/3/5, 3/4/1, 4/6/14}
    {\brick{\lo}{\hi}{1}{\v}}
  \foreach \lo/\hi/\v in {1/2/1, 2/3/4, 4/5/5, 5/6/9}
    {\brick{\lo}{\hi}{0}{\v}}
  \foreach \i in {0,...,5}
    \node[note] at (\i, -0.9) {\i};
  \begin{scope}[xshift=55mm]
    \foreach \n/\lo/\hi/\d/\v in {n0/0/6/0/23, n1/0/3/1/8, n2/3/6/1/15,
        n3/0/1/2/3, n4/1/3/2/5, n5/3/4/2/1, n6/4/6/2/14, n7/1/2/3/1,
        n8/2/3/3/4, n9/4/5/3/5, n10/5/6/3/9}
      \node[tn] (\n) at ({9mm * ((\lo + \hi) / 2 - 0.5)}, {23.4mm - 10mm * \d}) {\v};
    \foreach \a/\b in {n0/n1, n0/n2, n1/n3, n1/n4, n2/n5, n2/n6, n4/n7,
        n4/n8, n6/n9, n6/n10}
      \draw[draw=ink!55,thin] (\a) -- (\b);
    \foreach \n/\i in {n3/0, n7/1, n8/2, n5/3, n9/4, n10/5}
      \node[note,anchor=north] at ([yshift=-0.5mm]\n.south) {\i};
  \end{scope}
\end{tikzpicture}
\end{center}
On the right is the same tree drawn the usual way, one circle per brick. When \(n\) is not a power of two, the two halves of a range can differ in size, so some leaves sit higher than others. Here \(n = 6\), and the leaves 0 and 3 stop one level early.

\subsubsection{Lazy Propagation}
Let's add 2 to the range \([2, 6)\). The update walks down the tree just like a query. A brick outside the range is skipped. A brick partly inside is split, and we visit both halves. A brick fully inside is updated at once using \textup{(U1)}. The walk stops there.

\begin{center}
\begin{tikzpicture}[x=6.5mm,y=6.5mm,plain/.style={},
    gone/.style={draw=muted,densely dashed,text=muted}]
  \brick{0}{6}{3}{23}
  \brick{0}{3}{2}{8}
  \brick[pb]{3}{6}{2}{15}
  \brick[pc]{0}{1}{1}{3}
  \brick{1}{3}{1}{5}
  \brick[gone]{3}{4}{1}{\textcolor{muted}{1}}
  \brick[gone]{4}{6}{1}{\textcolor{muted}{14}}
  \brick[pc]{1}{2}{0}{1}
  \brick[pb]{2}{3}{0}{4}
  \brick[gone]{4}{5}{0}{\textcolor{muted}{5}}
  \brick[gone]{5}{6}{0}{\textcolor{muted}{9}}
  \foreach \i in {0,...,5}
    \node[note] at (\i, -0.9) {\i};
  \draw[decorate,decoration={brace,mirror,amplitude=4pt},muted]
    (1.55, -1.45) -- (5.45, -1.45)
    node[midway,below=5pt,note] {add 2 to \([2, 6)\)};
  \begin{scope}[xshift=55mm]
    \foreach \n/\lo/\hi/\d/\v/\s in {n0/0/6/0/23/plain, n1/0/3/1/8/plain,
        n2/3/6/1/15/pb, n3/0/1/2/3/pc, n4/1/3/2/5/plain, n5/3/4/2/1/gone,
        n6/4/6/2/14/gone, n7/1/2/3/1/pc, n8/2/3/3/4/pb, n9/4/5/3/5/gone,
        n10/5/6/3/9/gone}
      \node[tn,\s] (\n) at ({9mm * ((\lo + \hi) / 2 - 0.5)}, {23.4mm - 10mm * \d}) {\v};
    \foreach \a/\b/\s in {n0/n1/plain, n0/n2/plain, n1/n3/plain, n1/n4/plain,
        n2/n5/gone, n2/n6/gone, n4/n7/plain, n4/n8/plain, n6/n9/gone,
        n6/n10/gone}
      \draw[draw=ink!55,thin,\s] (\a) -- (\b);
    \foreach \n/\i in {n3/0, n7/1, n8/2, n5/3, n9/4, n10/5}
      \node[note,anchor=north] at ([yshift=-0.5mm]\n.south) {\i};
  \end{scope}
\end{tikzpicture}
\end{center}
The amber bricks are fully inside, the grey ones are outside, and the dashed ones are never visited. The walk stops at \([3, 6)\) without going down, even though it covers three elements.

Each amber brick gets its new sum straight from its old one. The leaf \([2, 3)\) becomes \(4 + 2 = 6\), and \([3, 6)\) becomes \(15 + 2 \cdot 3 = 21\). On the way back up, every brick that was split recomputes its sum from its two halves. These are the blue bricks.

\begin{center}
\begin{tikzpicture}[x=6.5mm,y=6.5mm,plain/.style={},
    gone/.style={draw=muted,densely dashed,text=muted}]
  \brick[pa]{0}{6}{3}{31}
  \brick[pa]{0}{3}{2}{10}
  \brick[pb]{3}{6}{2}{21}
  \node[font=\scriptsize,text=warncolor,anchor=east] at (5.4, 2.4) {\(+2\)};
  \brick{0}{1}{1}{3}
  \brick[pa]{1}{3}{1}{7}
  \brick[gone]{3}{4}{1}{\textcolor{muted}{1}}
  \brick[gone]{4}{6}{1}{\textcolor{muted}{14}}
  \brick{1}{2}{0}{1}
  \brick[pb]{2}{3}{0}{6}
  \brick[gone]{4}{5}{0}{\textcolor{muted}{5}}
  \brick[gone]{5}{6}{0}{\textcolor{muted}{9}}
  \foreach \i in {0,...,5}
    \node[note] at (\i, -0.9) {\i};
  \begin{scope}[xshift=55mm]
    \foreach \n/\lo/\hi/\d/\v/\s in {n0/0/6/0/31/pa, n1/0/3/1/10/pa,
        n2/3/6/1/21/pb, n3/0/1/2/3/plain, n4/1/3/2/7/pa, n5/3/4/2/1/gone,
        n6/4/6/2/14/gone, n7/1/2/3/1/plain, n8/2/3/3/6/pb, n9/4/5/3/5/gone,
        n10/5/6/3/9/gone}
      \node[tn,\s] (\n) at ({9mm * ((\lo + \hi) / 2 - 0.5)}, {23.4mm - 10mm * \d}) {\v};
    \node[font=\scriptsize,text=warncolor,anchor=west] at (n2.east) {\(+2\)};
    \foreach \a/\b/\s in {n0/n1/plain, n0/n2/plain, n1/n3/plain, n1/n4/plain,
        n2/n5/gone, n2/n6/gone, n4/n7/plain, n4/n8/plain, n6/n9/gone,
        n6/n10/gone}
      \draw[draw=ink!55,thin,\s] (\a) -- (\b);
    \foreach \n/\i in {n3/0, n7/1, n8/2, n5/3, n9/4, n10/5}
      \node[note,anchor=north] at ([yshift=-0.5mm]\n.south) {\i};
  \end{scope}
\end{tikzpicture}
\end{center}
The sums from \([3, 6)\) upward are right, but the dashed bricks below \([3, 6)\) still hold their old values. So \([3, 6)\) keeps the update as a tag, the red \(+2\), which means every element below it still needs 2 added. The tag waits there until a later walk has to go below it.

The idea of lazy propagation is that we are allowed to be lazy. We do not update the leaves below a node until someone actually asks about them. The nodes above never look at single leaves anyway. They only need the value of a whole range, and \textup{(U1)} gives that in one step.

\subsubsection{Merging Tags}
Before any query, let's add 2 to \([3, 6)\) once more. The walk goes straight to \([3, 6)\), which is fully inside, and stops there. Its sum becomes \(21 + 2 \cdot 3 = 27\). But \([3, 6)\) still holds the first \(+2\), which has not reached its children yet. By \textup{(U2)} the two updates merge into one, so the brick keeps a single tag \(+4\) to be pushed later.

\begin{center}
\begin{tikzpicture}[x=6.5mm,y=6.5mm,plain/.style={},
    gone/.style={draw=muted,densely dashed,text=muted}]
  \brick[pa]{0}{6}{3}{37}
  \brick{0}{3}{2}{10}
  \brick[pb]{3}{6}{2}{27}
  \node[font=\scriptsize,text=warncolor,anchor=east] at (5.4, 2.4) {\(+4\)};
  \brick{0}{1}{1}{3}
  \brick{1}{3}{1}{7}
  \brick[gone]{3}{4}{1}{\textcolor{muted}{1}}
  \brick[gone]{4}{6}{1}{\textcolor{muted}{14}}
  \brick{1}{2}{0}{1}
  \brick{2}{3}{0}{6}
  \brick[gone]{4}{5}{0}{\textcolor{muted}{5}}
  \brick[gone]{5}{6}{0}{\textcolor{muted}{9}}
  \foreach \i in {0,...,5}
    \node[note] at (\i, -0.9) {\i};
  \begin{scope}[xshift=55mm]
    \foreach \n/\lo/\hi/\d/\v/\s in {n0/0/6/0/37/pa, n1/0/3/1/10/plain,
        n2/3/6/1/27/pb, n3/0/1/2/3/plain, n4/1/3/2/7/plain, n5/3/4/2/1/gone,
        n6/4/6/2/14/gone, n7/1/2/3/1/plain, n8/2/3/3/6/plain, n9/4/5/3/5/gone,
        n10/5/6/3/9/gone}
      \node[tn,\s] (\n) at ({9mm * ((\lo + \hi) / 2 - 0.5)}, {23.4mm - 10mm * \d}) {\v};
    \node[font=\scriptsize,text=warncolor,anchor=west] at (n2.east) {\(+4\)};
    \foreach \a/\b/\s in {n0/n1/plain, n0/n2/plain, n1/n3/plain, n1/n4/plain,
        n2/n5/gone, n2/n6/gone, n4/n7/plain, n4/n8/plain, n6/n9/gone,
        n6/n10/gone}
      \draw[draw=ink!55,thin,\s] (\a) -- (\b);
    \foreach \n/\i in {n3/0, n7/1, n8/2, n5/3, n9/4, n10/5}
      \node[note,anchor=north] at ([yshift=-0.5mm]\n.south) {\i};
  \end{scope}
\end{tikzpicture}
\end{center}
The bricks below \([3, 6)\) have now missed two updates, yet they still wait for just one tag.

\subsubsection{Push}
Now let's ask for the element at position 4, the range \([4, 5)\). The walk reaches \([3, 6)\). It is only partly inside, so the walk has to go down. But the bricks below are out of date. So before going down, \([3, 6)\) pushes its tag to its two halves. Each half applies the \(+4\) to its own sum using \textup{(U1)} and keeps the tag for its own children. The tag on \([3, 6)\) is then reset to \(\mathrm{id}\), because nothing below it is waiting anymore.

\begin{center}
\begin{tikzpicture}[x=6.5mm,y=6.5mm,plain/.style={},
    gone/.style={draw=muted,densely dashed,text=muted}]
  \brick{0}{6}{3}{37}
  \brick{0}{3}{2}{10}
  \brick[hot]{3}{6}{2}{27}
  \brick{0}{1}{1}{3}
  \brick{1}{3}{1}{7}
  \brick[pb]{3}{4}{1}{5}
  \brick[pb]{4}{6}{1}{}
  \node[font=\small] at (4.25, 1.2) {22};
  \node[font=\scriptsize,text=warncolor,anchor=east] at (5.4, 1.2) {\(+4\)};
  \brick{1}{2}{0}{1}
  \brick{2}{3}{0}{6}
  \brick[gone]{4}{5}{0}{\textcolor{muted}{5}}
  \brick[gone]{5}{6}{0}{\textcolor{muted}{9}}
  \foreach \i in {0,...,5}
    \node[note] at (\i, -0.9) {\i};
  \begin{scope}[xshift=55mm]
    \foreach \n/\lo/\hi/\d/\v/\s in {n0/0/6/0/37/plain, n1/0/3/1/10/plain,
        n2/3/6/1/27/hot, n3/0/1/2/3/plain, n4/1/3/2/7/plain, n5/3/4/2/5/pb,
        n6/4/6/2/22/pb, n7/1/2/3/1/plain, n8/2/3/3/6/plain, n9/4/5/3/5/gone,
        n10/5/6/3/9/gone}
      \node[tn,\s] (\n) at ({9mm * ((\lo + \hi) / 2 - 0.5)}, {23.4mm - 10mm * \d}) {\v};
    \node[font=\scriptsize,text=warncolor,anchor=west] at (n6.east) {\(+4\)};
    \foreach \a/\b/\s in {n0/n1/plain, n0/n2/plain, n1/n3/plain, n1/n4/plain,
        n2/n5/plain, n2/n6/plain, n4/n7/plain, n4/n8/plain, n6/n9/gone,
        n6/n10/gone}
      \draw[draw=ink!55,thin,\s] (\a) -- (\b);
    \foreach \n/\i in {n3/0, n7/1, n8/2, n5/3, n9/4, n10/5}
      \node[note,anchor=north] at ([yshift=-0.5mm]\n.south) {\i};
  \end{scope}
\end{tikzpicture}
\end{center}
The red brick is the one that pushed. The leaf 3 became \(1 + 4 = 5\), and \([4, 6)\) became \(14 + 4 \cdot 2 = 22\) and now holds the \(+4\) itself. Both updates travel down together in one push.

The same happens one level down. \([4, 6)\) is partly inside too, so it pushes its \(+4\) to the leaves 4 and 5 before the walk goes there.\enlargethispage{2\baselineskip}

\begin{center}
\begin{tikzpicture}[x=6.5mm,y=6.5mm,plain/.style={}]
  \brick{0}{6}{3}{37}
  \brick{0}{3}{2}{10}
  \brick{3}{6}{2}{27}
  \brick{0}{1}{1}{3}
  \brick{1}{3}{1}{7}
  \brick{3}{4}{1}{5}
  \brick[hot]{4}{6}{1}{22}
  \brick{1}{2}{0}{1}
  \brick{2}{3}{0}{6}
  \brick[pb]{4}{5}{0}{9}
  \brick[pb]{5}{6}{0}{13}
  \foreach \i in {0,...,5}
    \node[note] at (\i, -0.9) {\i};
  \draw[decorate,decoration={brace,mirror,amplitude=4pt},muted]
    (3.55, -1.45) -- (4.45, -1.45)
    node[midway,below=5pt,note] {the answer, 9};
  \begin{scope}[xshift=55mm]
    \foreach \n/\lo/\hi/\d/\v/\s in {n0/0/6/0/37/plain, n1/0/3/1/10/plain,
        n2/3/6/1/27/plain, n3/0/1/2/3/plain, n4/1/3/2/7/plain, n5/3/4/2/5/plain,
        n6/4/6/2/22/hot, n7/1/2/3/1/plain, n8/2/3/3/6/plain, n9/4/5/3/9/pb,
        n10/5/6/3/13/pb}
      \node[tn,\s] (\n) at ({9mm * ((\lo + \hi) / 2 - 0.5)}, {23.4mm - 10mm * \d}) {\v};
    \foreach \a/\b in {n0/n1, n0/n2, n1/n3, n1/n4, n2/n5, n2/n6, n4/n7,
        n4/n8, n6/n9, n6/n10}
      \draw[draw=ink!55,thin] (\a) -- (\b);
    \foreach \n/\i in {n3/0, n7/1, n8/2, n5/3, n9/4, n10/5}
      \node[note,anchor=north] at ([yshift=-0.5mm]\n.south) {\i};
  \end{scope}
\end{tikzpicture}
\end{center}
The walk ends at the leaf \([4, 5)\), which now holds the right value 9. A tag only moves when a walk passes through its brick, one level at a time.

\subsubsection{The Six Lines}
Every rule from the green box is one line of code. Here are the first eight lines of the notebook code, set up for adding with sum.

\excerpt[SegTree,merge,apply,comp]{../notebook/code/ds/segtree.cpp}{1}{8}

\begin{notes}[Reading the code]
\begin{lines}
\item[Lines 2--3] \code|T| is the type of a value (sum). \code|L| is the type of a tag (the amount to add).
\item[Line 4] \code|e| is the identity of the query from \textup{(Q2)}. An empty range has sum 0.
\item[Line 5] \code|id| is the update that changes nothing from \textup{(U3)}, which is adding 0.
\item[Line 6] \code|merge| is the query \(x \oplus y\) from \textup{(Q1)}.
\item[Line 7] \code|apply| is \(x \circ f\) on a whole range from \textup{(U1)}. The old sum alone is not enough. Adding \(f\) to \(k\) elements raises the sum by \(f \cdot k\), so the code also passes the length of the range.
\item[Line 8] \code|comp(f, g)| merges two tags from \textup{(U2)}. It means \(g\) first and then \(f\).
\end{lines}
\end{notes}

To use the tree for another update or query, only these lines change. Here are two more setups.

\textbf{Adding with sum and minimum together.} A value can hold several answers at once. Here it is a struct with the sum and the minimum of its range, and \code|merge| combines both. \\
Adding \(f\) raises the sum by \(f\) times the length and the minimum by \(f\).

\begin{cpp}
  struct T { ll sum, mn; };
  typedef ll L;
  T e = {0, LLONG_MAX};
  L id = 0;
  T merge(T a, T b) {
    return {a.sum + b.sum, min(a.mn, b.mn)};
  }
  T apply(T x, L f, int len) {
    return {x.sum + f * len, x.mn + f};
  }
  L comp(L f, L g) { return f + g; }
\end{cpp}

\textbf{Multiplying and adding with sum.} A tag is a pair \((a, b)\) for \(x \mapsto ax + b\). On a whole range it turns the sum into \(a\) times the sum plus \(b\) times the length. Merging uses the formula from the grey box. In a real problem every product is taken modulo the given prime.

\begin{cpp}
  typedef ll T;
  typedef pair<ll, ll> L; // x -> a x + b
  T e = 0;
  L id = {1, 0};
  T merge(T a, T b) { return a + b; }
  T apply(T x, L f, int len) {
    return f.first * x + f.second * len;
  }
  L comp(L f, L g) {
    return {f.first * g.first, f.first * g.second + f.second};
  }
\end{cpp}

\textbf{Any other pair.} To build the lines for a new pair, answer three questions. The query side only needs \code|e|, the result of an empty range.
\begin{enumerate}[label=\textup{(U\arabic*)},leftmargin=*]
\item Given the value of a range and its length, what is its value after \(f\)? This is \code|apply|.
\item Doing \(g\) and then \(f\), which single update is that? This is \code|comp|.
\item Which update changes nothing? This is \code|id|.
\end{enumerate}

\begin{center}
\begin{tabular}{lllll}
\toprule
Update & Query & \code|id| & \code|apply(x, f, len)| & \code|comp(f, g)| \\
\midrule
add \(f\) & sum & \(0\) & \(x + f \cdot \mathit{len}\) & \(f + g\) \\
add \(f\) & minimum & \(0\) & \(x + f\) & \(f + g\) \\
set to \(f\) & sum & none & \(f \cdot \mathit{len}\) & \(f\) \\
set to \(f\) & minimum & none & \(f\) & \(f\) \\
multiply by \(f\) & sum & \(1\) & \(f \cdot x\) & \(f \cdot g\) \\
maximum with \(f\) & maximum & \(-\infty\) & \(\max(x, f)\) & \(\max(f, g)\) \\
flip a 0/1 array & count of ones & \(0\) & \(\mathit{len} - x\) if \(f = 1\) & \(f \operatorname{xor} g\) \\
\bottomrule
\end{tabular}
\end{center}

The \code|id| does not have to be a real update. Its only job is to say that nothing is waiting to be pushed. Setting has no update that changes nothing, so none is any value that never appears as a real tag, such as \code|LLONG_MIN|. With none, \code|apply| returns \(x\) and \code|comp| returns \(g\).

With some creativity and careful coding, many more operations fit into these six lines.

\subsubsection{The Code}
Brace yourself. Understand the concepts first, and every line of code below should look familiar.

\excerpt[SegTree,build,merge]{../notebook/code/ds/segtree.cpp}{10}{22}

\begin{notes}[Reading the code]
\begin{lines}
\item[Lines 10--12] The length \code|n|, the values \code|t| and the tags \code|lz|, one value and one tag per node.
\item[Line 13] Node 1 is the root, and the halves of node \code|v| are \code|2 * v| and \code|2 * v + 1|. \\
Four times \code|n| nodes are always enough. \\
This builds an empty tree with \code|n| elements, every value \code|e| and every tag \code|id|.
\item[Lines 14--16] Builds the tree from an array. Starts at the root, which covers \code|[0, n)|.
\item[Line 17] \code|build| fills node \code|v|, which covers \code|[lo, hi)|.
\item[Line 18] A leaf is a brick of width one. \\
Its range \code|[lo, hi)| holds a single element, so \code|lo| is its index in the array. \\
It copies its element from the array and stops.
\item[Line 19] Every other node splits at \code|m|, just like the bricks.
\item[Line 20] It builds its left half \code|[lo, m)| and its right half \code|[m, hi)| first.
\item[Line 21] Then it merges the two halves into its own value.
\end{lines}
\end{notes}

The two helpers are the two figures from Lazy Propagation and Push.

\excerpt[put,push,apply,comp]{../notebook/code/ds/segtree.cpp}{23}{29}

\begin{notes}[Reading the code]
\begin{lines}
\item[Line 24] \code|put| updates a whole brick. \\
Its value takes \code|f| at once by \textup{(U1)}. \\
Its tag merges \code|f| on top of the waiting one by \textup{(U2)}.
\item[Lines 27--28] \code|push| hands the tag to both halves with \code|put|, then resets it to \code|id| from \textup{(U3)}.
\end{lines}
\end{notes}

The update is the walk from Lazy Propagation.

\excerpt[update,put,push,merge]{../notebook/code/ds/segtree.cpp}{31}{40}

\begin{notes}[Reading the code]
\begin{lines}
\item[Line 33] A brick outside the range is skipped.
\item[Line 34] A brick fully inside takes the update with \code|put|, and the walk stops.
\item[Line 36] A brick partly inside pushes its tag before the walk goes down.
\item[Line 39] On the way back up, the brick recomputes its value from its halves.
\end{lines}
\end{notes}

The query walks the same way.

\excerpt[query,push,merge]{../notebook/code/ds/segtree.cpp}{42}{50}

\begin{notes}[Reading the code]
\begin{lines}
\item[Line 44] A brick outside gives \code|e|, which changes nothing in a merge by \textup{(Q2)}.
\item[Line 45] A brick fully inside gives its value.
\item[Line 47] A brick partly inside pushes first, exactly as in Push.
\item[Lines 48--49] The left half is merged first, so the elements stay in order.
\end{lines}
\end{notes}

Last, \code|set| replaces a single element.

\excerpt[set,push,merge]{../notebook/code/ds/segtree.cpp}{52}{60}

\begin{notes}[Reading the code]
\begin{lines}
\item[Line 54] At the leaf, the new value replaces the old one.
\item[Line 56] Push on the way down. \\
Otherwise a tag waiting above would later change the new value.
\item[Line 59] Recompute on the way back up.
\end{lines}
\end{notes}

\subsubsection{Practice}
After fully understanding the lazy segment tree, it is time to put it to the test.

\begin{center}
\begin{tabular}{@{}lll@{}}
\toprule
Problem & Level & Practices \\
\midrule
\href{https://cses.fi/problemset/task/1651}{Range Update Queries} & Easy & adding, then reading one element \\
\href{https://cses.fi/problemset/task/1735}{Range Updates and Sums} & Medium & setting and adding mixed together \\
\href{https://atcoder.jp/contests/practice2/tasks/practice2_k}{Range Affine Range Sum} & Medium & multiplying and adding \\
\href{https://atcoder.jp/contests/practice2/tasks/practice2_l}{Lazy Segment Tree} & Medium & a value with several fields \\
\href{https://atcoder.jp/contests/abc357/tasks/abc357_f}{Two Sequence Queries} & Medium & a sum of products \\
\href{https://atcoder.jp/contests/abc322/tasks/abc322_f}{Vacation Query} & Medium & longest run of ones, with flips \\
\href{https://codeforces.com/problemset/problem/1477/B}{Nezzar and Binary String} & Medium & setting, solved backwards \\
\href{https://cses.fi/problemset/task/1736}{Polynomial Queries} & Hard & a tag that depends on position \\
\href{https://codeforces.com/problemset/problem/242/E}{XOR on Segment} & Hard & xor on a range \\
\href{https://codeforces.com/problemset/problem/1320/C}{World of Darkraft: Battle for Azathoth} & Hard & a sweep with adding and max \\
\href{https://codeforces.com/problemset/problem/1555/E}{Boring Segments} & Hard & two pointers with adding and min \\
\href{https://codeforces.com/problemset/problem/1263/E}{Editor} & Hard & bracket depth, prefix min and max \\
\href{https://codeforces.com/problemset/problem/620/E}{New Year Tree} & Hard & colors as a bitmask on a subtree \\
\href{https://codeforces.com/problemset/problem/52/C}{Circular RMQ} & Hard & ranges that wrap around \\
\href{https://codeforces.com/problemset/problem/1179/C}{Serge and Dining Room} & Hard & walking down to a position \\
\href{https://codeforces.com/problemset/problem/1557/D}{Ezzat and Grid} & Hard & a dp with max over ranges \\
\href{https://codeforces.com/problemset/problem/718/C}{Sasha and Array} & Very hard & matrices as values and tags \\
\href{https://codeforces.com/problemset/problem/558/E}{A Simple Task} & Very hard & sorting a substring \\
\href{https://codeforces.com/problemset/problem/803/G}{Periodic RMQ Problem} & Very hard & a huge array, compressed \\
\href{https://codeforces.com/problemset/problem/817/F}{MEX Queries} & Very hard & set, clear and flip, then the first gap \\
\href{https://codeforces.com/problemset/problem/1114/F}{Please, another Queries on Array?} & Very hard & products and a bitmask of primes \\
\href{https://codeforces.com/problemset/problem/446/C}{DZY Loves Fibonacci Numbers} & Very hard & Fibonacci sequences as tags \\
\href{https://codeforces.com/problemset/problem/1638/E}{Colorful Operations} & Very hard & coloring ranges, adding per color \\
\href{https://codeforces.com/problemset/problem/580/E}{Kefa and Watch} & Very hard & setting, with string hashes \\
\href{https://codeforces.com/problemset/problem/266/E}{More Queries to Array...} & Very hard & setting, with sums of powers \\
\href{https://codeforces.com/problemset/problem/1373/G}{Pawns} & Very hard & adding, with a clever value \\
\href{https://codeforces.com/problemset/problem/1371/F}{Raging Thunder} & Very hard & flipping directions, many fields \\
\href{https://codeforces.com/problemset/problem/280/D}{k-Maximum Subsequence Sum} & Very hard & negating a range, best subarray \\
\bottomrule
\end{tabular}
\end{center}

\topic{Implicit Treap}[\log n]

\begin{problem}
Given an array \(a\) of length \(n\), we can do the following in \(\bigO{\log n}\).
\begin{itemize}
\item Insert or erase an element at any position.
\item Split the array into several pieces, then join the pieces back in any order.
\end{itemize}

\begin{center}
\begin{tikzpicture}[x=6.5mm,y=6.5mm]
  \foreach \c [count=\i from 0] in {A,B,C}
    \node[cell,pa] at (\i, 0) {\c};
  \foreach \c [count=\i from 3] in {D,E}
    \node[cell,pb] at (\i, 0) {\c};
  \draw[cut] (2.5, 0.8) -- (2.5, -0.8);
  \node[note,text=warncolor] at (2.5, 1.2) {split};
  \foreach \c [count=\i from 0] in {A,B,C}
    \node[cell,pa] at (\i - 0.4, -2.2) {\c};
  \foreach \c [count=\i from 3] in {D,E}
    \node[cell,pb] at (\i + 0.4, -2.2) {\c};
  \node[note] at (0.6, -3.1) {left piece};
  \node[note] at (3.9, -3.1) {right piece};

  \begin{scope}[xshift=55mm]
    \foreach \c [count=\i from 0] in {A,B}
      \node[cell,pa] at (\i, 0) {\c};
    \foreach \c [count=\i from 2] in {C,D,E}
      \node[cell,pb] at (\i, 0) {\c};
    \node[cell,pc] at (5, 0) {F};
    \draw[cut] (1.5, 0.8) -- (1.5, -0.8);
    \draw[cut] (4.5, 0.8) -- (4.5, -0.8);
    \node[note,text=warncolor] at (3, 1.2) {multiple splits};
    \foreach \c [count=\i from 0] in {A,B}
      \node[cell,pa] at (\i - 0.4, -2.2) {\c};
    \foreach \c [count=\i from 2] in {C,D,E}
      \node[cell,pb] at (\i, -2.2) {\c};
    \node[cell,pc] at (5.4, -2.2) {F};
    \node[note] at (0.1, -3.1) {\([0, l)\)};
    \node[note] at (3, -3.1) {\([l, r)\)};
    \node[note] at (5.4, -3.1) {\([r, n)\)};
    \foreach \c [count=\i from 0] in {C,D,E}
      \node[cell,pb] at (\i, -4.6) {\c};
    \foreach \c [count=\i from 3] in {A,B}
      \node[cell,pa] at (\i, -4.6) {\c};
    \node[cell,pc] at (5, -4.6) {F};
    \node[note,text=defcolor] at (2.5, -5.5) {join in a new order};
  \end{scope}
\end{tikzpicture}
\end{center}

\begin{itemize}
\item Run every range update and range query a lazy segment tree supports.
\end{itemize}
\end{problem}

\subsubsection{Tree Structure}
Every node holds one element. Reading the tree in order gives the array. In order means the left subtree first, then the node itself, then the right subtree.

\begin{center}
\begin{tikzpicture}[tree]
  \node[tn] {B}
    child {node {A}}
    child {node {D} child {node {C}} child {node {E}}};
  \begin{scope}[xshift=48mm,yshift=-11mm,x=6.5mm,y=6.5mm]
    \foreach \c [count=\i from 0] in {A,B,C,D,E} {
      \node[cell] at (\i, 0) {\c};
      \node[note] at (\i, -0.9) {\i};
    }
  \end{scope}
\end{tikzpicture}
\end{center}

\subsubsection{Subtrees are blocks}
The nodes of a subtree are read one after another, so every subtree is a block of consecutive elements. A node is therefore one element and also one block, the block of its subtree. This is what lets a single node store a summary of a whole block, such as its sum, and it is why one node can stand for a whole range.

\begin{center}
\begin{tikzpicture}[tree]
  \node[tn] {B}
    child {node[pa] {A}}
    child {node[pb] {D} child {node[pb] {C}} child {node[pb] {E}}};
  \begin{scope}[xshift=48mm,yshift=-11mm,x=6.5mm,y=6.5mm]
    \node[cell,pa] at (0, 0) {A};
    \node[cell] at (1, 0) {B};
    \foreach \c [count=\i from 2] in {C,D,E}
      \node[cell,pb] at (\i, 0) {\c};
  \end{scope}
\end{tikzpicture}
\end{center}
The subtree of \texttt{D} is the block \texttt{C D E}, and the subtree of \texttt{A} is the block \texttt{A}.

\subsubsection{Implicit positions}
No node stores its position. Instead every node stores the size of its subtree, which is the length of its block. Inside that block the left subtree comes first, so the node sits at position \(L\), the size of its left subtree. That size is just the number stored at the left child, or \(0\) if there is none. To find position \(k\), start at the root. If \(k < L\), go left and look for \(k\). If \(k = L\), the root is the answer. If \(k > L\), go right and look for \(k - L - 1\), skipping the left block and the root.

\begin{center}
\begin{tikzpicture}[tree]
  \node[tn,label={[note]right:5}] {B}
    child {node[label={[note]left:1}] {A}}
    child {node[label={[note]right:3}] {D}
      child {node[label={[note]left:1}] {C}}
      child {node[label={[note]right:1}] {E}}};
  \begin{scope}[xshift=48mm,yshift=-11mm,x=6.5mm,y=6.5mm]
    \foreach \c [count=\i from 0] in {A,B,C,D,E} {
      \node[cell] at (\i, 0) {\c};
      \node[note] at (\i, -0.9) {\i};
    }
  \end{scope}
\end{tikzpicture}
\end{center}
The small numbers are subtree sizes. Let's look for position 2.

Start at the root. \texttt{B} has \(L = 1\) and \(2 > 1\), so \texttt{A} and \texttt{B} come before position 2. We move right to \texttt{D} and look for \(2 - 1 - 1 = 0\) inside its block.

\begin{center}
\begin{tikzpicture}[tree]
  \node[tn,pc,label={[note]right:5}] {B}
    child {node[pc,label={[note]left:1}] {A}}
    child {node[hot,label={[note]right:3}] {D}
      child {node[label={[note]left:1}] {C}}
      child {node[label={[note]right:1}] {E}}};
  \begin{scope}[xshift=48mm,yshift=-11mm,x=6.5mm,y=6.5mm]
    \foreach \c [count=\i from 0] in {A,B,C,D,E} {
      \node[cell] at (\i, 0) {\c};
      \node[note] at (\i, -0.9) {\i};
    }
    \node[cell,pc] at (0, 0) {A};
    \node[cell,pc] at (1, 0) {B};
  \end{scope}
\end{tikzpicture}
\end{center}

At \texttt{D}, \(L = 1\) and \(0 < 1\), so the answer lies in the left block, and \texttt{D} and \texttt{E} come after it. We move left to \texttt{C} and look for 0.

\begin{center}
\begin{tikzpicture}[tree]
  \node[tn,pc,label={[note]right:5}] {B}
    child {node[pc,label={[note]left:1}] {A}}
    child {node[pc,label={[note]right:3}] {D}
      child {node[hot,label={[note]left:1}] {C}}
      child {node[pc,label={[note]right:1}] {E}}};
  \begin{scope}[xshift=48mm,yshift=-11mm,x=6.5mm,y=6.5mm]
    \foreach \c [count=\i from 0] in {A,B,C,D,E} {
      \node[cell] at (\i, 0) {\c};
      \node[note] at (\i, -0.9) {\i};
    }
    \foreach \c/\i in {A/0,B/1,D/3,E/4}
      \node[cell,pc] at (\i, 0) {\c};
    \node[cell,hot] at (2, 0) {C};
  \end{scope}
\end{tikzpicture}
\end{center}

At \texttt{C}, \(L = 0\) and \(0 = 0\), so \texttt{C} is the answer, and it indeed sits at position 2.

\subsubsection{Priorities}
Many trees spell the same array, so we need a rule that picks one. Every node gets a random priority, and every parent must have a larger priority than its children. The root is then the node with the largest priority, the nodes before it form its left subtree, the nodes after it form its right subtree, and the same rule repeats inside each subtree. So the priorities decide the shape completely.

\begin{center}
\begin{tikzpicture}[tree]
  \node[tn,label={[note]right:9}] at (0, 0) {B}
    child {node[label={[note]left:3}] {A}}
    child {node[label={[note]right:7}] {D}
      child {node[label={[note]left:5}] {C}}
      child {node[label={[note]right:2}] {E}}};
  \begin{scope}[xshift=-9.5mm,yshift=-46mm,x=6.5mm,y=6.5mm]
    \foreach \c/\p [count=\i from 0] in {A/3,B/9,C/5,D/7,E/2} {
      \node[cell] at (\i, 0) {\c};
      \node[note] at (\i, -0.9) {\p};
    }
  \end{scope}
  \node[tn,label={[note]right:9}] at (75mm, 0) {E}
    child {node[label={[note]left:8}] {A}
      child[missing]
      child {node[label={[note]right:6}] {C}
        child {node[label={[note]left:2}] {B}}
        child {node[label={[note]right:4}] {D}}}}
    child[missing];
  \begin{scope}[xshift=56mm,yshift=-46mm,x=6.5mm,y=6.5mm]
    \foreach \c/\p [count=\i from 0] in {A/8,B/2,C/6,D/4,E/9} {
      \node[cell] at (\i, 0) {\c};
      \node[note] at (\i, -0.9) {\p};
    }
  \end{scope}
\end{tikzpicture}
\end{center}
The small numbers are priorities. On the left, \texttt{B} has the largest, so it is the root, and among \texttt{C D E} the largest is \texttt{D}. On the right, the same array gets other priorities. Now \texttt{E} is the root and the tree is deeper, yet it still reads \texttt{A B C D E}.

Random priorities put the root at a random position, like a random pivot in quicksort. So the expected depth is \(\bigO{\log n}\). \textit{That's just how randomness works.} Every walk from the root to a node is therefore short. Priorities only affect speed. Every operation stays correct whatever shape the tree has.

\subsubsection{Split}
We can split a sequence after its first \(k\) elements. This gives two blocks, which means two trees, and their shapes can differ from the original tree. We will walk through this on a longer sequence.

\begin{center}
\begin{tikzpicture}[tree,
    level 1/.style={sibling distance=44mm},
    level 2/.style={sibling distance=18mm},
    level 3/.style={sibling distance=12mm}]
  \node[tn,label={[note]right:9}] {E}
    child {node[label={[note]left:7}] {B}
      child {node[label={[note]left:3}] {A}}
      child {node[label={[note]right:6}] {D}
        child {node[label={[note]left:2}] {C}}
        child[missing]}}
    child {node[label={[note]right:8}] {G}
      child {node[label={[note]left:4}] {F}}
      child {node[label={[note]right:5}] {H}}};
  \begin{scope}[xshift=50mm,yshift=-16mm,x=6.5mm,y=6.5mm]
    \foreach \c/\p [count=\i from 0] in {A/3,B/7,C/2,D/6,E/9,F/4,G/8,H/5} {
      \node[cell] at (\i, 0) {\c};
      \node[note] at (\i, -0.9) {\p};
    }
  \end{scope}
\end{tikzpicture}
\end{center}
The small numbers are priorities. Let's split after three elements.

Start at the root. Four elements come before \texttt{E}, so \texttt{E} sits at position 4, while the left tree only takes positions 0, 1 and 2. So \texttt{E} lies after the cut, and \texttt{E} with its whole right subtree belongs to the right tree. We move to \texttt{B}.

\begin{center}
\begin{tikzpicture}[tree,
    level 1/.style={sibling distance=44mm},
    level 2/.style={sibling distance=18mm},
    level 3/.style={sibling distance=12mm}]
  \node[tn,pb,label={[note]right:9}] {E}
    child {node[hot,label={[note]left:7}] {B}
      child {node[label={[note]left:3}] {A}}
      child {node[label={[note]right:6}] {D}
        child {node[label={[note]left:2}] {C}}
        child[missing]}
      edge from parent[draw=none]}
    child {node[pb,label={[note]right:8}] {G}
      child {node[pb,label={[note]left:4}] {F}}
      child {node[pb,label={[note]right:5}] {H}}};
  \begin{scope}[xshift=50mm,yshift=-16mm,x=6.5mm,y=6.5mm]
    \foreach \c/\p [count=\i from 0] in {A/3,B/7,C/2,D/6} {
      \node[cell] at (\i, 0) {\c};
      \node[note] at (\i, -0.9) {\p};
    }
    \foreach \c/\p [count=\i from 4] in {E/9,F/4,G/8,H/5} {
      \node[cell,pb] at (\i, 0) {\c};
      \node[note] at (\i, -0.9) {\p};
    }
    \draw[cut] (2.5, 0.8) -- (2.5, -0.8);
  \end{scope}
\end{tikzpicture}
\end{center}

At \texttt{B}, one element comes before it, so \texttt{B} sits at position 1, inside the left tree. So \texttt{B} and its whole left subtree belong to the left tree. The cut must still pass through its right subtree after \(3 - 1 - 1 = 1\) more element. We move to \texttt{D}.

\begin{center}
\begin{tikzpicture}[tree,
    level 1/.style={sibling distance=44mm},
    level 2/.style={sibling distance=18mm},
    level 3/.style={sibling distance=12mm}]
  \node[tn,pb,label={[note]right:9}] {E}
    child {node[pa,label={[note]left:7}] {B}
      child {node[pa,label={[note]left:3}] {A}}
      child {node[hot,label={[note]right:6}] {D}
        child {node[label={[note]left:2}] {C}}
        child[missing]
        edge from parent[draw=none]}
      edge from parent[draw=none]}
    child {node[pb,label={[note]right:8}] {G}
      child {node[pb,label={[note]left:4}] {F}}
      child {node[pb,label={[note]right:5}] {H}}};
  \begin{scope}[xshift=50mm,yshift=-16mm,x=6.5mm,y=6.5mm]
    \foreach \c/\p [count=\i from 0] in {A/3,B/7} {
      \node[cell,pa] at (\i, 0) {\c};
      \node[note] at (\i, -0.9) {\p};
    }
    \foreach \c/\p [count=\i from 2] in {C/2,D/6} {
      \node[cell] at (\i, 0) {\c};
      \node[note] at (\i, -0.9) {\p};
    }
    \foreach \c/\p [count=\i from 4] in {E/9,F/4,G/8,H/5} {
      \node[cell,pb] at (\i, 0) {\c};
      \node[note] at (\i, -0.9) {\p};
    }
    \draw[cut] (2.5, 0.8) -- (2.5, -0.8);
  \end{scope}
\end{tikzpicture}
\end{center}

At \texttt{D}, one element is still to be cut, and one element, \texttt{C}, comes before \texttt{D} inside its block. So \texttt{D} lies after the cut, and \texttt{D} with its right subtree belongs to the right tree. We move to \texttt{C}.

\begin{center}
\begin{tikzpicture}[tree,
    level 1/.style={sibling distance=44mm},
    level 2/.style={sibling distance=18mm},
    level 3/.style={sibling distance=12mm}]
  \node[tn,pb,label={[note]right:9}] {E}
    child {node[pa,label={[note]left:7}] {B}
      child {node[pa,label={[note]left:3}] {A}}
      child {node[pb,label={[note]right:6}] {D}
        child {node[hot,label={[note]left:2}] {C}
          edge from parent[draw=none]}
        child[missing]
        edge from parent[draw=none]}
      edge from parent[draw=none]}
    child {node[pb,label={[note]right:8}] {G}
      child {node[pb,label={[note]left:4}] {F}}
      child {node[pb,label={[note]right:5}] {H}}};
  \begin{scope}[xshift=50mm,yshift=-16mm,x=6.5mm,y=6.5mm]
    \foreach \c/\p [count=\i from 0] in {A/3,B/7} {
      \node[cell,pa] at (\i, 0) {\c};
      \node[note] at (\i, -0.9) {\p};
    }
    \node[cell] at (2, 0) {C};
    \node[note] at (2, -0.9) {2};
    \foreach \c/\p [count=\i from 3] in {D/6,E/9,F/4,G/8,H/5} {
      \node[cell,pb] at (\i, 0) {\c};
      \node[note] at (\i, -0.9) {\p};
    }
    \draw[cut] (2.5, 0.8) -- (2.5, -0.8);
  \end{scope}
\end{tikzpicture}
\end{center}

At \texttt{C}, nothing comes before it and one element is still to be cut. So \texttt{C} lies before the cut and belongs to the left tree. The walk would go on into its right subtree, but that subtree is empty, so the walk ends here.

Every node now has its color, but the pieces are still loose. Split finishes on the way back up the same path, starting from the bottom. Each node receives two pieces from the call below it, keeps the piece on its own side of the cut and passes the other one up. The table lists the two trees each node returns, written as the elements they read, with the root in bold.

\begin{center}
\begin{tabular}{@{}lll@{}}
\toprule
Node & Left tree & Right tree \\
\midrule
\texttt{C} & \texttt{\textbf{C}} & \(\varnothing\) \\
\texttt{D} & \texttt{\textbf{C}} & \texttt{\textbf{D}} \\
\texttt{B} & \texttt{A \textbf{B} C} & \texttt{\textbf{D}} \\
\texttt{E} & \texttt{A \textbf{B} C} & \texttt{D \textbf{E} F G H} \\
\bottomrule
\end{tabular}
\end{center}

The dotted lines are the two new edges.

\begin{center}
\begin{tikzpicture}[tree,
    level 1/.style={sibling distance=44mm},
    level 2/.style={sibling distance=18mm},
    level 3/.style={sibling distance=12mm}]
  \node[tn,pb,label={[note]right:9}] (e) {E}
    child {node[pa,label={[note]left:7}] (b) {B}
      child {node[pa,label={[note]left:3}] {A}}
      child {node[pb,label={[note]right:6}] (d) {D}
        child {node[pa,label={[note]left:2}] (c) {C}
          edge from parent[draw=none]}
        child[missing]
        edge from parent[draw=none]}
      edge from parent[draw=none]}
    child {node[pb,label={[note]right:8}] {G}
      child {node[pb,label={[note]left:4}] {F}}
      child {node[pb,label={[note]right:5}] {H}}};
  \draw[defcolor,thick,densely dotted] (c) -- (b);
  \draw[defcolor,thick,densely dotted] (d) -- (e);
  \begin{scope}[xshift=50mm,yshift=-16mm,x=6.5mm,y=6.5mm]
    \foreach \c/\p [count=\i from 0] in {A/3,B/7,C/2} {
      \node[cell,pa] at (\i, 0) {\c};
      \node[note] at (\i, -0.9) {\p};
    }
    \foreach \c/\p [count=\i from 3] in {D/6,E/9,F/4,G/8,H/5} {
      \node[cell,pb] at (\i, 0) {\c};
      \node[note] at (\i, -0.9) {\p};
    }
    \draw[cut] (2.5, 0.8) -- (2.5, -0.8);
  \end{scope}
\end{tikzpicture}
\end{center}

The last row of the table is what the whole call returns. The left tree reads \texttt{A B C} and the right tree reads \texttt{D E F G H}. Both have shapes the original tree never had.

\begin{center}
\begin{tikzpicture}[tree,
    level 1/.style={sibling distance=18mm},
    level 2/.style={sibling distance=18mm}]
  \node[tn,pa,label={[note]right:7}] at (0, 0) {B}
    child {node[pa,label={[note]left:3}] {A}}
    child {node[pa,label={[note]right:2}] {C}};
  \node[tn,pb,label={[note]right:9}] at (62mm, 0) {E}
    child {node[pb,label={[note]left:6}] {D}}
    child {node[pb,label={[note]right:8}] {G}
      child {node[pb,label={[note]left:4}] {F}}
      child {node[pb,label={[note]right:5}] {H}}};
  \begin{scope}[xshift=-6.5mm,yshift=-38mm,x=6.5mm,y=6.5mm]
    \foreach \c [count=\i from 0] in {A,B,C}
      \node[cell,pa] at (\i, 0) {\c};
    \node[note] at (1, -0.9) {left tree};
  \end{scope}
  \begin{scope}[xshift=49mm,yshift=-38mm,x=6.5mm,y=6.5mm]
    \foreach \c [count=\i from 0] in {D,E,F,G,H}
      \node[cell,pb] at (\i, 0) {\c};
    \node[note] at (2, -0.9) {right tree};
  \end{scope}
\end{tikzpicture}
\end{center}

\subsubsection{Join}
We can join two trees \code|a| and \code|b| into one, with every element of \code|a| before every element of \code|b|.

Join compares the two roots and puts the one with the larger priority on top. The interesting part is what the winner keeps. Suppose the root of \code|a| wins. Its left subtree comes right before it, and nothing from \code|b| can come before it, because all of \code|b| comes after all of \code|a|. So the left subtree stays as its left child. After the root come its right subtree and then all of \code|b|, in that order. Gluing these two is another join, one level lower, and the result becomes the new right child.

When the root of \code|b| wins, it is the mirror image. Its right subtree comes right after it, and nothing from \code|a| can come after it. So the right subtree stays as its right child. Before the root come all of \code|a| and then its left subtree. Gluing these two is another join, and the result becomes the new left child. When one of the trees is empty, the result is simply the other tree.

\begin{center}
\begin{tikzpicture}[blk/.style={draw=ink!55,minimum height=6.5mm,
    inner sep=0pt,font=\footnotesize,anchor=west}]
  \node[note,anchor=west] at (0, 7mm) {the root of \texttt{a} wins};
  \node[blk,pa,minimum width=28mm] at (0, 0) {left subtree};
  \node[blk,pa,hot,minimum width=12mm] at (28mm, 0) {root};
  \node[blk,pa,minimum width=28mm] at (40mm, 0) {right subtree};
  \node[blk,pb,minimum width=32mm] at (68mm, 0) {all of \texttt{b}};
  \draw[decorate,decoration={brace,mirror,amplitude=4pt},muted]
    (0.5mm, -4.5mm) -- (27.5mm, -4.5mm)
    node[midway,below=5pt,note] {stays as the left child};
  \draw[decorate,decoration={brace,mirror,amplitude=4pt},muted]
    (40.5mm, -4.5mm) -- (99.5mm, -4.5mm)
    node[midway,below=5pt,note] {joined next, becomes the right child};
  \begin{scope}[yshift=-28mm]
    \node[note,anchor=west] at (0, 7mm) {the root of \texttt{b} wins};
    \node[blk,pa,minimum width=32mm] at (0, 0) {all of \texttt{a}};
    \node[blk,pb,minimum width=28mm] at (32mm, 0) {left subtree};
    \node[blk,pb,hot,minimum width=12mm] at (60mm, 0) {root};
    \node[blk,pb,minimum width=28mm] at (72mm, 0) {right subtree};
    \draw[decorate,decoration={brace,mirror,amplitude=4pt},muted]
      (0.5mm, -4.5mm) -- (59.5mm, -4.5mm)
      node[midway,below=5pt,note] {joined next, becomes the left child};
    \draw[decorate,decoration={brace,mirror,amplitude=4pt},muted]
      (72.5mm, -4.5mm) -- (99.5mm, -4.5mm)
      node[midway,below=5pt,note] {stays as the right child};
  \end{scope}
\end{tikzpicture}
\end{center}

So the priorities only decide which node goes higher. Which subtree the winner keeps is what puts every element in its place, and that depends only on which tree the winner comes from.

Let's join the two trees we just got from split. We should get the original tree back.

Compare the two roots. \texttt{E} has priority 9 and \texttt{B} has 7, so \texttt{E} goes on top. Everything in the right tree comes after the left tree, so \texttt{E} keeps its right subtree. Its left side must hold the whole left tree followed by its old left subtree \texttt{D}, so we join those two next.

\begin{center}
\begin{tikzpicture}[tree,
    level 1/.style={sibling distance=44mm},
    level 2/.style={sibling distance=18mm},
    level 3/.style={sibling distance=12mm}]
  \node[tn,pb,label={[note]right:9}] {E}
    child {node[slot] {?}}
    child {node[pb,label={[note]right:8}] {G}
      child {node[pb,label={[note]left:4}] {F}}
      child {node[pb,label={[note]right:5}] {H}}};
  \node[note] at (0, -40mm) {built so far};
  \begin{scope}[xshift=70mm,level 1/.style={sibling distance=16mm}]
    \node[tn,pa,hot,label={[note]right:7}] {B}
      child {node[pa,label={[note]left:3}] {A}}
      child {node[pa,label={[note]right:2}] {C}};
    \node[tn,pb,hot,label={[note]right:6}] at (28mm, 0) {D};
    \node[note] at (14mm, -40mm) {still to join};
  \end{scope}
\end{tikzpicture}
\end{center}

Compare \texttt{B} with 7 and \texttt{D} with 6. \texttt{B} goes on top and fills the open spot as the left child of \texttt{E}. \texttt{B} comes before \texttt{D}, so it keeps its left subtree \texttt{A}. Its right side must hold its old right subtree \texttt{C} followed by \texttt{D}, so we join those two next.

\begin{center}
\begin{tikzpicture}[tree,
    level 1/.style={sibling distance=44mm},
    level 2/.style={sibling distance=18mm},
    level 3/.style={sibling distance=12mm}]
  \node[tn,pb,label={[note]right:9}] {E}
    child {node[pa,label={[note]left:7}] {B}
      child {node[pa,label={[note]left:3}] {A}}
      child {node[slot] {?}}
      edge from parent[new]}
    child {node[pb,label={[note]right:8}] {G}
      child {node[pb,label={[note]left:4}] {F}}
      child {node[pb,label={[note]right:5}] {H}}};
  \node[note] at (0, -40mm) {built so far};
  \begin{scope}[xshift=70mm]
    \node[tn,pa,hot,label={[note]right:2}] {C};
    \node[tn,pb,hot,label={[note]right:6}] at (28mm, 0) {D};
    \node[note] at (14mm, -40mm) {still to join};
  \end{scope}
\end{tikzpicture}
\end{center}

Compare \texttt{C} with 2 and \texttt{D} with 6. \texttt{D} goes on top and fills the open spot as the right child of \texttt{B}. \texttt{D} comes after \texttt{C}, so it keeps its right subtree, which is empty. Its left side must hold \texttt{C} followed by its old left subtree, which is also empty.

\begin{center}
\begin{tikzpicture}[tree,
    level 1/.style={sibling distance=44mm},
    level 2/.style={sibling distance=18mm},
    level 3/.style={sibling distance=12mm}]
  \node[tn,pb,label={[note]right:9}] {E}
    child {node[pa,label={[note]left:7}] {B}
      child {node[pa,label={[note]left:3}] {A}}
      child {node[pb,label={[note]right:6}] {D}
        child {node[slot] {?}}
        child[missing]
        edge from parent[new]}
      edge from parent[new]}
    child {node[pb,label={[note]right:8}] {G}
      child {node[pb,label={[note]left:4}] {F}}
      child {node[pb,label={[note]right:5}] {H}}};
  \node[note] at (0, -40mm) {built so far};
  \begin{scope}[xshift=70mm]
    \node[tn,pa,hot,label={[note]right:2}] {C};
    \node[note] at (28mm, 0) {\(\varnothing\)};
    \node[note] at (14mm, -40mm) {still to join};
  \end{scope}
\end{tikzpicture}
\end{center}

Joining \texttt{C} with an empty tree just gives \texttt{C}. So \texttt{C} fills the last open spot as the left child of \texttt{D}. The join is done, and we got the original tree back. The dotted lines are the three edges join created.

\begin{center}
\begin{tikzpicture}[tree,
    level 1/.style={sibling distance=44mm},
    level 2/.style={sibling distance=18mm},
    level 3/.style={sibling distance=12mm}]
  \node[tn,pb,label={[note]right:9}] {E}
    child {node[pa,label={[note]left:7}] {B}
      child {node[pa,label={[note]left:3}] {A}}
      child {node[pb,label={[note]right:6}] {D}
        child {node[pa,label={[note]left:2}] {C}
          edge from parent[new]}
        child[missing]
        edge from parent[new]}
      edge from parent[new]}
    child {node[pb,label={[note]right:8}] {G}
      child {node[pb,label={[note]left:4}] {F}}
      child {node[pb,label={[note]right:5}] {H}}};
  \begin{scope}[xshift=50mm,yshift=-16mm,x=6.5mm,y=6.5mm]
    \foreach \c/\p [count=\i from 0] in {A/3,B/7,C/2} {
      \node[cell,pa] at (\i, 0) {\c};
      \node[note] at (\i, -0.9) {\p};
    }
    \foreach \c/\p [count=\i from 3] in {D/6,E/9,F/4,G/8,H/5} {
      \node[cell,pb] at (\i, 0) {\c};
      \node[note] at (\i, -0.9) {\p};
    }
  \end{scope}
\end{tikzpicture}
\end{center}

Join only compares priorities and never looks at sizes, because the order is already fixed by which tree each element comes from.

\subsubsection{Operations}
With split and join, every operation in the green box is a short recipe. We cut the sequence into pieces, work on the pieces and join them back. Each recipe uses a fixed number of splits and joins, so every operation costs \(\bigO{\log n}\).

\textbf{Insert} a value \(x\) at position \(i\). Split after \(i\) elements, then join the left piece, a new tree holding only \(x\) and the right piece, in that order.

\textbf{Erase} the element at position \(i\). Split after \(i\) elements, then split the right piece after one more element. The middle piece holds only the element to erase. Join the two outer pieces and drop the middle one.

\begin{center}
\begin{tikzpicture}[x=6.5mm,y=6.5mm,px/.style={fill=defcolor!22}]
  \foreach \c [count=\i from 0] in {A,B,C,D,E,F,G,H}
    \node[cell] at (\i, 0) {\c};
  \draw[cut] (2.5, 0.8) -- (2.5, -0.8);
  \node[note,text=warncolor] at (2.5, 1.2) {split after 3};
  \foreach \c [count=\i from 0] in {A,B,C}
    \node[cell,pa] at (\i - 0.5, -2.2) {\c};
  \node[cell,px] at (3, -2.2) {X};
  \foreach \c [count=\i from 3] in {D,E,F,G,H}
    \node[cell,pb] at (\i + 1.5, -2.2) {\c};
  \node[note] at (3, -3.1) {new tree};
  \foreach \c [count=\i from 0] in {A,B,C}
    \node[cell,pa] at (\i, -4.4) {\c};
  \node[cell,px] at (3, -4.4) {X};
  \foreach \c [count=\i from 4] in {D,E,F,G,H}
    \node[cell,pb] at (\i, -4.4) {\c};
  \node[note,text=defcolor] at (4, -5.3) {join the three pieces};

  \begin{scope}[xshift=72mm]
    \foreach \c [count=\i from 0] in {A,B,C,D,E,F,G,H}
      \node[cell] at (\i, 0) {\c};
    \draw[cut] (2.5, 0.8) -- (2.5, -0.8);
    \draw[cut] (3.5, 0.8) -- (3.5, -0.8);
    \node[note,text=warncolor] at (3, 1.2) {split twice};
    \foreach \c [count=\i from 0] in {A,B,C}
      \node[cell,pa] at (\i - 0.5, -2.2) {\c};
    \node[cell,pc] at (3, -2.2) {D};
    \foreach \c [count=\i from 4] in {E,F,G,H}
      \node[cell,pb] at (\i + 0.5, -2.2) {\c};
    \node[note] at (3, -3.1) {dropped};
    \foreach \c [count=\i from 0] in {A,B,C}
      \node[cell,pa] at (\i, -4.4) {\c};
    \foreach \c [count=\i from 3] in {E,F,G,H}
      \node[cell,pb] at (\i, -4.4) {\c};
    \node[note,text=defcolor] at (3, -5.3) {join the outer pieces};
  \end{scope}
\end{tikzpicture}
\end{center}

\textbf{Work on a block} \(a[l, r)\). Split after \(l\) elements, then split the right piece after \(r - l\) more. The middle piece is now a whole tree, and its root stands for exactly the block. A range update, a range query or a reversal acts on that one root, as the next part explains. Then join the three pieces back in their old order.

\textbf{Rearrange blocks.} Cut the sequence into as many pieces as needed and join them in a new order. Moving a block, swapping two blocks and cyclically shifting a block all work this way.

\begin{center}
\begin{tikzpicture}[x=6.5mm,y=6.5mm,px/.style={fill=defcolor!22}]
  \foreach \c [count=\i from 0] in {A,B,C,D,E,F,G,H}
    \node[cell] at (\i, 0) {\c};
  \draw[cut] (1.5, 0.8) -- (1.5, -0.8);
  \draw[cut] (4.5, 0.8) -- (4.5, -0.8);
  \node[note,text=warncolor] at (3, 1.2) {cut out \([2, 5)\)};
  \foreach \c [count=\i from 0] in {A,B}
    \node[cell,pa] at (\i - 0.5, -2.2) {\c};
  \foreach \c [count=\i from 2] in {C,D,E}
    \node[cell,pb] at (\i, -2.2) {\c};
  \foreach \c [count=\i from 5] in {F,G,H}
    \node[cell,pc] at (\i + 0.5, -2.2) {\c};
  \node[note] at (3, -3.1) {one tree, act on its root};
  \foreach \c [count=\i from 0] in {A,B}
    \node[cell,pa] at (\i, -4.4) {\c};
  \foreach \c [count=\i from 2] in {E,D,C}
    \node[cell,pb] at (\i, -4.4) {\c};
  \foreach \c [count=\i from 5] in {F,G,H}
    \node[cell,pc] at (\i, -4.4) {\c};
  \node[note,text=defcolor] at (3.5, -5.3) {reversed, joined back};

  \begin{scope}[xshift=72mm]
    \foreach \c [count=\i from 0] in {A,B,C,D,E,F,G,H}
      \node[cell] at (\i, 0) {\c};
    \draw[cut] (1.5, 0.8) -- (1.5, -0.8);
    \draw[cut] (4.5, 0.8) -- (4.5, -0.8);
    \draw[cut] (5.5, 0.8) -- (5.5, -0.8);
    \node[note,text=warncolor] at (3.5, 1.2) {cut into four};
    \foreach \c [count=\i from 0] in {A,B}
      \node[cell,pa] at (\i - 0.6, -2.2) {\c};
    \foreach \c [count=\i from 2] in {C,D,E}
      \node[cell,pb] at (\i - 0.2, -2.2) {\c};
    \node[cell,px] at (5.2, -2.2) {F};
    \foreach \c [count=\i from 6] in {G,H}
      \node[cell,pc] at (\i + 0.6, -2.2) {\c};
    \foreach \c [count=\i from 0] in {A,B}
      \node[cell,pa] at (\i, -4.4) {\c};
    \node[cell,px] at (2, -4.4) {F};
    \foreach \c [count=\i from 3] in {C,D,E}
      \node[cell,pb] at (\i, -4.4) {\c};
    \foreach \c [count=\i from 6] in {G,H}
      \node[cell,pc] at (\i, -4.4) {\c};
    \node[note,text=defcolor] at (3.5, -5.3) {joined in a new order};
  \end{scope}
\end{tikzpicture}
\end{center}
On the right, the block \texttt{C D E F} is shifted right by one, so \texttt{F} wraps around to its front.

\subsubsection{Lazy tags}
Besides its element, every node stores the sum of its block. Let \texttt{A} to \texttt{H} hold the values 1 to 8, and let's add 10 to the block \texttt{A B C}.

We cut the block out, which gives the left tree from split, and update only its root. The value of \texttt{B} grows by 10. The sum of its block grows by 10 for each of its 3 elements, from 6 to 36. The children \texttt{A} and \texttt{C} are not touched. Instead \texttt{B} keeps a note saying that its children still owe 10. This note is called a lazy tag.

\begin{center}
\begin{tikzpicture}[tree,level 1/.style={sibling distance=22mm}]
  \node[tn,pa,label={[note]right:12},
    label={[font=\footnotesize,text=warncolor]above:\(+10\)}] {B}
    child {node[pa,label={[note]left:1}] {A}}
    child {node[pa,label={[note]right:3}] {C}};
  \node[note] at (0, -20mm) {block sum at \texttt{B} is \(36\)};
\end{tikzpicture}
\end{center}
The small numbers are values, and the red number is the tag.

A tag is delivered by pushing. Pushing a node applies its tag to both children the same way, adds it to their own tags and clears it at the node. Split and join push a node before they touch its children. One reason is reading them, because their values may still be stale. The other is moving them, because a tag is meant only for the children the node has when the tag is written.

Now we join \texttt{A B C} back. Join reaches \texttt{B} and is about to give it a new right child, so it pushes \texttt{B} first. \texttt{A} becomes 11, \texttt{C} becomes 13 and the tag at \texttt{B} is cleared. Only then does \texttt{C} move under \texttt{D}, with its update already inside.

\begin{center}
\begin{tikzpicture}[tree,
    level 1/.style={sibling distance=44mm},
    level 2/.style={sibling distance=18mm},
    level 3/.style={sibling distance=12mm}]
  \node[tn,pb,label={[note]right:5}] {E}
    child {node[pa,label={[note]left:12}] {B}
      child {node[pa,label={[note]left:11}] {A}}
      child {node[pb,label={[note]right:4}] {D}
        child {node[pa,label={[note]left:13}] {C}}
        child[missing]}}
    child {node[pb,label={[note]right:7}] {G}
      child {node[pb,label={[note]left:6}] {F}}
      child {node[pb,label={[note]right:8}] {H}}};
\end{tikzpicture}
\end{center}

Without that push, \texttt{C} would move under \texttt{D} still holding 3, and the tag would stay at \texttt{B}. The children of \texttt{B} are now \texttt{A} and \texttt{D}, so the tag would later add 10 to the whole block of \texttt{D}. But \texttt{D} was never part of the block we updated.

\begin{center}
\begin{tikzpicture}[tree,
    level 1/.style={sibling distance=44mm},
    level 2/.style={sibling distance=18mm},
    level 3/.style={sibling distance=12mm}]
  \node[tn,pb,label={[note]right:5}] {E}
    child {node[pa,label={[note]left:12},
        label={[font=\footnotesize,text=warncolor]above:\(+10\)}] {B}
      child {node[pa,label={[note]left:1}] {A}}
      child {node[pb,hot,label={[note]right:4}] {D}
        child {node[pa,label={[note]left:3}] {C}}
        child[missing]}}
    child {node[pb,label={[note]right:7}] {G}
      child {node[pb,label={[note]left:6}] {F}}
      child {node[pb,label={[note]right:8}] {H}}};
\end{tikzpicture}
\end{center}
This is the tree without the push. The tag at \texttt{B} now hangs above \texttt{D}, outlined in red.

\subsubsection{Reverse}
To reverse a block, we mirror its tree, swapping the left and right child at every node.

\begin{center}
\begin{tikzpicture}[tree,
    level 1/.style={sibling distance=20mm},
    level 2/.style={sibling distance=12mm}]
  \node[tn] at (0, 0) {B}
    child {node {A}}
    child {node {D} child {node {C}} child {node {E}}};
  \begin{scope}[xshift=-13mm,yshift=-34mm,x=6.5mm,y=6.5mm]
    \foreach \c [count=\i from 0] in {A,B,C,D,E}
      \node[cell] at (\i, 0) {\c};
  \end{scope}
  \draw[->,muted,thick] (28mm, -10mm) -- (44mm, -10mm)
    node[midway,above,note] {mirror};
  \node[tn] at (72mm, 0) {B}
    child {node {D} child {node {E}} child {node {C}}}
    child {node {A}};
  \begin{scope}[xshift=59mm,yshift=-34mm,x=6.5mm,y=6.5mm]
    \foreach \c [count=\i from 0] in {E,D,C,B,A}
      \node[cell] at (\i, 0) {\c};
  \end{scope}
\end{tikzpicture}
\end{center}

Reading a tree always gives the left part, then the root, then the right part. After mirroring the old right part comes first and the old left part comes last, and both are mirrored as well. So the whole block reads backwards.

Mirroring every node takes as long as the block is big, so the treap does it lazily. A flag on a node says that its subtree should be mirrored but has not been yet. Reversing a block cuts it out and flips the flag on its root. Pushing a flagged node does one level of the work. It swaps the two children of the node and flips the flag on each of them, because their subtrees still need mirroring inside.

Let's reverse the whole tree and then look for position 0. Reversing only puts a flag on \texttt{B}. To find position 0 we walk down and push on the way.

\begin{center}
\begin{tikzpicture}[tree,
    level 1/.style={sibling distance=16mm},
    level 2/.style={sibling distance=10mm},
    flag/.style={label={[font=\footnotesize,text=warncolor]above:rev}}]
  \node[tn,flag] at (0, 0) {B}
    child {node {A}}
    child {node {D} child {node {C}} child {node {E}}};
  \draw[->,muted,thick] (16mm, -8mm) -- (28mm, -8mm)
    node[midway,above,note] {push \texttt{B}};
  \node[tn] at (46mm, 0) {B}
    child {node[flag] {D} child {node {C}} child {node {E}}}
    child {node[flag] {A}};
  \draw[->,muted,thick] (62mm, -8mm) -- (74mm, -8mm)
    node[midway,above,note] {push \texttt{D}};
  \node[tn] at (92mm, 0) {B}
    child {node {D} child {node[flag] {E}} child {node[flag] {C}}}
    child {node[flag] {A}};
\end{tikzpicture}
\end{center}
The red labels are flags. Pushing \texttt{B} swaps its children, so its left part is now the block of \texttt{D} with 3 elements, and position 0 lies inside it. Pushing \texttt{D} swaps its children too, and position 0 turns out to be \texttt{E}, the first element of \texttt{E D C B A}. The flags at \texttt{A}, \texttt{C} and \texttt{E} are still waiting because nobody has walked there yet.

This walk also shows why split pushes a node before it reads the size of the left subtree. The left part of \texttt{B} had 1 element before the push and 3 after it.

The flag is flipped, not set, so that two reversals cancel each other. A reversal never changes the sum of a block, so the block sums stay correct without extra work. The same holds for the minimum or the maximum, but not for anything that depends on order.

\subsubsection{The code}
Every idea above is a few lines of the notebook code. The excerpts keep the line numbers of the notebook file.

The first eight lines are the six rules from the segment tree, with the same contract. Here they describe range add with range sum.

\excerpt[Treap,merge,apply,comp]{../notebook/code/ds/treap.cpp}{1}{8}

The nodes and the storage come next.

\excerpt[make]{../notebook/code/ds/treap.cpp}{9}{18}

\begin{notes}[Reading the code]
\begin{itemize}
\item Line 10. Children \code|l| and \code|r|, subtree size \code|sz|, priority \code|pri|, own value \code|v|, block sum \code|s|, tag \code|lz| and reverse flag \code|rv|.
\item Line 12. All nodes live in one vector and point at each other by index. Index 0 is the empty tree, with size 0 and sum \code|e|, so a missing child is just 0.
\item Line 13. A fixed seed is fine in ICPC, which has no hacking, and it makes every run reproducible.
\item Line 14. The tree that holds the whole sequence, empty at the start.
\item Lines 15--18. A tree of one node with a random priority. It returns the index of the node.
\end{itemize}
\end{notes}

When the children of a node change, \code|pull| recomputes its size and its block sum, reading the left child, the node and the right child in that order.

\excerpt[pull,merge]{../notebook/code/ds/treap.cpp}{29}{33}

Split follows the walk we drew, one call per node on the path.

\excerpt[split,push,pull]{../notebook/code/ds/treap.cpp}{34}{46}

\begin{notes}[Reading the code]
\begin{itemize}
\item Line 35. Splitting an empty tree gives two empty trees.
\item Line 36. Push before reading the children, because a waiting flag may swap them.
\item Line 37. \code|ls| is \(L\), the number of elements before the root inside its tree.
\item Lines 38--42. The cut falls before the root. The left subtree is split. The root keeps the part after the cut as its new left child and passes the part before the cut up.
\item Lines 43--45. The cut falls after the root. The right subtree is split after the remaining \code|k - ls - 1| elements. The root keeps the part before the cut as its new right child and passes the part after the cut up.
\end{itemize}
\end{notes}

Join compares the two roots and keeps one subtree of the winner, exactly as in the block picture.

\excerpt[join,push,pull]{../notebook/code/ds/treap.cpp}{47}{55}

\begin{notes}[Reading the code]
\begin{itemize}
\item Line 48. If one tree is empty, the result is the other. One of the two indices is 0, so \code|a ^ b| is the other one.
\item Lines 49--52. The root of \code|a| wins, keeps its left subtree, and its right subtree is joined with \code|b|.
\item Lines 53--54. The root of \code|b| wins, keeps its right subtree, and \code|a| is joined with its left subtree.
\item Both cases push before changing a child and pull after.
\end{itemize}
\end{notes}

The operations are the recipes from before.

\excerpt[insert,erase,on,update,reverse,query,split,join,make,put]{../notebook/code/ds/treap.cpp}{57}{80}

\begin{notes}[Reading the code]
\begin{itemize}
\item Lines 57--60. Split after \code|i|, then join the left piece, the new node and the right piece.
\item Lines 61--64. Split after \code|i| and after one more, then join the outer pieces. The single node in the middle is never used again.
\item Lines 66--71. \code|on| cuts out \([l, r)\), lets \code|f| act on the root of the middle tree and joins the three pieces back. The template lets \code|f| be any lambda.
\item Lines 72--80. \code|update| tags the root, \code|reverse| flips its flag and \code|query| reads its block sum.
\end{itemize}
\end{notes}

Finally the lazy part, which split and join have been calling all along.

\excerpt[put,push,apply,comp]{../notebook/code/ds/treap.cpp}{19}{28}

\begin{notes}[Reading the code]
\begin{itemize}
\item Line 20. The empty node must stay empty.
\item Line 21. The own value counts as one element, and the new tag goes on top of the waiting one.
\item Line 22. The block sum counts \code|sz| elements.
\item Line 25. A reference to the node, safe because nothing here adds nodes. A reference kept across \code|make| could point at memory the vector has already left.
\item Line 26. A waiting reversal swaps the children and passes the flag to both.
\item Line 27. The waiting tag goes to both children, and both notes are cleared.
\end{itemize}
\end{notes}

\begin{pitfall}
Never keep a \code|Node&| across a call to \code|make|, never read the children of a node before pushing it, and never reverse with a merge that depends on order. When the positions never change, use a segment tree instead, which is several times faster.
\end{pitfall}

\topic{Persistent Segment Tree}[\log n]

\begin{problem}
Given an array \(a\) of length \(n\), we can do the following in \(\bigO{\log n}\).
\begin{itemize}
\item Update one element or query a range, like a segment tree.
\item Every update makes a new version of the array and keeps the old one.\\
An update can start from any old version.
\end{itemize}

\begin{center}
\begin{tikzpicture}[x=6.5mm,y=6.5mm]
  \foreach \v [count=\i from 0] in {0,0,0,0}
    \node[cell] at (\i, 0) {\v};
  \node[note,anchor=east] at (-0.8, 0) {version 0};
  \foreach \v [count=\i from 0] in {0,3,0,0}
    \node[cell] at (\i, -2.4) {\v};
  \node[cell,pb] at (1, -2.4) {3};
  \node[note,anchor=east] at (-0.8, -2.4) {version 1};
  \draw[->,muted,thick] (1.5, -0.6) -- (1.5, -1.8)
    node[midway,right,note] {\texttt{a[1] += 3}};
  \foreach \v [count=\i from 0] in {0,3,0,2}
    \node[cell] at (\i - 3.5, -4.8) {\v};
  \node[cell,pb] at (-0.5, -4.8) {2};
  \node[note] at (-2, -5.7) {version 2};
  \foreach \v [count=\i from 0] in {4,3,0,0}
    \node[cell] at (\i + 3.5, -4.8) {\v};
  \node[cell,pb] at (3.5, -4.8) {4};
  \node[note] at (5, -5.7) {version 3};
  \draw[->,muted,thick] (0.9, -3) -- (-1.6, -4.2)
    node[midway,above left,note] {\texttt{a[3] += 2}};
  \draw[->,muted,thick] (2.1, -3) -- (4.6, -4.2)
    node[midway,above right,note] {\texttt{a[0] += 4}};
\end{tikzpicture}
\end{center}
\end{problem}

\subsubsection{A segment tree forgets}
To add \(x\) to \(a_i\), a segment tree walks from the root down to the leaf of \(i\) and adds \(x\) to the sum at every node on the way. Those are exactly the nodes whose blocks contain \(i\). Each new sum is written over the old one, so once the update is done the old array is gone for good.

\begin{center}
\begin{tikzpicture}[tree,level 1/.style={sibling distance=28mm},
    gone/.style={draw=muted,densely dashed,text=muted}]
  \begin{scope}[edge from parent/.style={draw=muted,densely dashed}]
    \node[tn,gone] at (0, 0) {0}
      child {node[gone] {0}
        child {node[gone,label={[note]below:0}] {0}}
        child {node[gone,label={[note]below:1}] {0}}}
      child {node[gone] {0}
        child {node[gone,label={[note]below:2}] {0}}
        child {node[gone,label={[note]below:3}] {0}}};
  \end{scope}
  \node[note] at (0, -34mm) {before, now gone};
  \draw[->,muted,thick] (30mm, -11mm) -- (44mm, -11mm)
    node[midway,above,note] {\texttt{a[1] += 3}};
  \node[tn,hot] at (74mm, 0) {3}
    child {node[hot] {3}
      child {node[label={[note]below:0}] {0}}
      child {node[hot,label={[note]below:1}] {3}}}
    child {node {0}
      child {node[label={[note]below:2}] {0}}
      child {node[label={[note]below:3}] {0}}};
  \node[note] at (74mm, -34mm) {after};
\end{tikzpicture}
\end{center}
The red nodes are the path. Their old sums are overwritten and lost.

\subsubsection{Copy the path}
The obvious way to keep the old array is to copy all of it before every update, which costs \(\bigO{n}\) each time. A persistent segment tree copies only the path. It never writes over a node. Instead it makes a copy of every node on the path and puts the new sum in the copy. Each copy keeps its old child that is off the path, so everything off the path is shared by the two versions instead of copied.

\begin{center}
\begin{tikzpicture}[tree,level 1/.style={sibling distance=28mm}]
  \node[tn] at (0, 0) {0}
    child {node {0}
      child {node[label={[note]below:0}] {0}}
      child {node[label={[note]below:1}] {0}}}
    child {node {0}
      child {node[label={[note]below:2}] {0}}
      child {node[label={[note]below:3}] {0}}};
  \node[note] at (0, -34mm) {version 0, still there};
  \draw[->,muted,thick] (30mm, -11mm) -- (44mm, -11mm)
    node[midway,above,note] {\texttt{a[1] += 3}};
  \node[tn,pb] at (74mm, 0) {3}
    child {node[pb] {3}
      child {node[label={[note]below:0}] {0}}
      child {node[pb,label={[note]below:1}] {3}}}
    child {node {0}
      child {node[label={[note]below:2}] {0}}
      child {node[label={[note]below:3}] {0}}};
  \node[note] at (74mm, -34mm) {version 1};
\end{tikzpicture}
\end{center}
The amber nodes are the copies. The white nodes of version~1 are not copies at all. They are the very same nodes that version~0 uses.

Version~0 starts at the old root and meets only old nodes, so it reads exactly as before. Version~1 starts at the new root. So a version is nothing more than a root, and an update costs one path, which is \(\bigO{\log n}\) new nodes.

\subsubsection{Storing the nodes}
All nodes live in one big array \code|t|, in the order they were made. A node stores its sum and the indices of its two children. So a child is just a number, and two parents share a child by storing the same number. An update appends the copies of its path to the end of \code|t| and returns the index of the new root. That index is the whole version, so we keep the roots in a list where \code|root[v]| is the root of version~\(v\).

So far version~0 was a full tree of zeros, built before the first update. The code never builds it. Instead \code|t[0]| is a node with sum~0 whose two children are node~0 itself. Walking down from node~0 never meets anything but zeros, however deep it goes. So node~0 acts like a tree of zeros of any size, and version~0 is simply root~0. The first update copies node~0 along its path, and every child it does not replace is still node~0.

\begin{center}
\begin{tikzpicture}[tree,level 1/.style={sibling distance=28mm},
    id/.style={label={[note,font=\ttfamily\footnotesize]#1}}]
  \node[tn,pc] (z) at (0, -11mm) {0};
  \draw[->,muted,thick] (z) to[out=250,in=190,looseness=12]
    node[pos=0.5,below left,note] {left} (z);
  \draw[->,muted,thick] (z) to[out=290,in=350,looseness=12]
    node[pos=0.5,below right,note] {right} (z);
  \node[note] at (0, -30mm) {node 0};
  \node[tn,pb,id={right:t[1]}] at (60mm, 0) {3}
    child {node[pb,id={left:t[2]}] {3}
      child {node[pc,id={left:t[0]}] {0}}
      child {node[pb,id={right:t[3]}] {3}}}
    child {node[pc,id={right:t[0]}] {0}};
  \node[note] at (60mm, -30mm) {version 1, built from root 0};
  \begin{scope}[x=6.5mm,y=6.5mm,xshift=-3.25mm,yshift=-54mm]
    \node[note] at (0.5, 1.8) {\texttt{root}};
    \foreach \v/\u/\c in {0/0/pc, 1/1/pb} {
      \node[note] at (\v, 0.9) {\v};
      \node[cell,\c] at (\v, 0) {\u};
    }
  \end{scope}
  \begin{scope}[x=6.5mm,y=6.5mm,xshift=50.25mm,yshift=-54mm]
    \node[note] at (1.5, 1.8) {\texttt{t}};
    \foreach \u/\l/\r/\s/\c in {0/0/0/0/pc, 1/2/0/3/pb, 2/0/3/3/pb,
        3/0/0/3/pb} {
      \node[note] at (\u, 0.9) {\u};
      \node[cell,\c] at (\u, 0) {\l};
      \node[cell,\c] at (\u, -1) {\r};
      \node[cell,\c] at (\u, -2) {\s};
    }
    \node[note,anchor=east] at (-0.6, 0) {\texttt{l}};
    \node[note,anchor=east] at (-0.6, -1) {\texttt{r}};
    \node[note,anchor=east] at (-0.6, -2) {\texttt{s}};
  \end{scope}
\end{tikzpicture}
\end{center}
The grey nodes are all the same node~0. Node~1 is the root of version~1. Its left child is node~2, and its right child is node~0, which stands for the block \([2, 4)\) full of zeros. So version~1 costs three new nodes and nothing else.

\subsubsection{Branching}
Since nothing is ever written over, every root ever made still works. An update can start from any of them, even an old one, and it simply builds another path. Here are the versions from the box. Every node is colored by the update that made it, amber for version~1, blue for version~2 and teal for version~3. Grey nodes are node~0.

\begin{center}
\begin{tikzpicture}[tree,
    level 1/.style={sibling distance=20mm},
    level 2/.style={sibling distance=10mm},
    pd/.style={fill=defcolor!18}]
  \node[tn,pb] at (0, 0) {3}
    child {node[pb] {3}
      child {node[pc,label={[note]below:0}] {0}}
      child {node[pb,label={[note]below:1}] {3}}}
    child {node[pc] {0}
      child {node[pc,label={[note]below:2}] {0}}
      child {node[pc,label={[note]below:3}] {0}}};
  \node[note] at (0, -34mm) {version 1};
  \node[tn,pa] at (50mm, 0) {5}
    child {node[pb] {3}
      child {node[pc,label={[note]below:0}] {0}}
      child {node[pb,label={[note]below:1}] {3}}}
    child {node[pa] {2}
      child {node[pc,label={[note]below:2}] {0}}
      child {node[pa,label={[note]below:3}] {2}}};
  \node[note] at (50mm, -34mm) {version 2, from version 1};
  \node[tn,pd] at (100mm, 0) {7}
    child {node[pd] {7}
      child {node[pd,label={[note]below:0}] {4}}
      child {node[pb,label={[note]below:1}] {3}}}
    child {node[pc] {0}
      child {node[pc,label={[note]below:2}] {0}}
      child {node[pc,label={[note]below:3}] {0}}};
  \node[note] at (100mm, -34mm) {version 3, from version 1};
\end{tikzpicture}
\end{center}
Versions~2 and~3 both grew from version~1, and both reuse its amber nodes. Still, neither can see what the other changed. The new nodes of version~2 hang below its own root only. No path from the root of version~3 ever reaches them.

In memory it is still one array. Each update appended its three nodes, and the list of roots grew by one.

\begin{center}
\begin{tikzpicture}[x=6.5mm,y=6.5mm,pd/.style={fill=defcolor!18}]
  \node[note] at (4.5, 1.8) {\texttt{t}};
  \foreach \u/\l/\r/\s/\c in {0/0/0/0/pc, 1/2/0/3/pb, 2/0/3/3/pb,
      3/0/0/3/pb, 4/2/5/5/pa, 5/0/6/2/pa, 6/0/0/2/pa, 7/8/0/7/pd,
      8/9/3/7/pd, 9/0/0/4/pd} {
    \node[note] at (\u, 0.9) {\u};
    \node[cell,\c] at (\u, 0) {\l};
    \node[cell,\c] at (\u, -1) {\r};
    \node[cell,\c] at (\u, -2) {\s};
  }
  \node[note,anchor=east] at (-0.6, 0) {\texttt{l}};
  \node[note,anchor=east] at (-0.6, -1) {\texttt{r}};
  \node[note,anchor=east] at (-0.6, -2) {\texttt{s}};
  \begin{scope}[xshift=13*6.5mm]
    \node[note] at (1.5, 1.8) {\texttt{root}};
    \foreach \v/\u/\c in {0/0/pc, 1/1/pb, 2/4/pa, 3/7/pd} {
      \node[note] at (\v, 0.9) {\v};
      \node[cell,\c] at (\v, 0) {\u};
    }
  \end{scope}
\end{tikzpicture}
\end{center}
Node~4, the root of version~2, keeps node~2 as its left child. Node~8 in version~3 keeps node~3 as its right child. Those are the shared amber nodes, and they are shared simply because two nodes store the same index.

\subsubsection{Query}
A query on version~\(v\) is the usual segment tree query, started at \code|root[v]|. It only reads nodes and never changes them, so every version can be queried at any time and in any order.

\subsubsection{The \(k\)th smallest}
The classic use of all this is finding the \(k\)th smallest element of a subarray \([l, r)\), with \(k\) counted from 0. Assume every value lies in \([0, m)\). Any array gets there by replacing each value with its rank among the distinct values, which keeps the order and makes \(m\) at most \(n\).

The trick is to build the tree over values instead of positions. Leaf \(x\) counts how many times \(x\) appears, so a node counts the elements whose values fall in its block. Start from version~0, where every count is zero, and add the elements one by one. Adding \(a_i\) adds 1 at leaf \(a_i\), so version~\(i\) counts the values in the prefix \([0, i)\).

\begin{center}
\begin{tikzpicture}[x=6.5mm,y=6.5mm]
  \foreach \v [count=\i from 0] in {3,0,2,3,1} {
    \node[note] at (\i, 0.9) {\i};
    \node[cell] at (\i, 0) {\v};
  }
  \node[note,anchor=east] at (-0.7, 0) {\(a\)};
  \begin{scope}[xshift=10*6.5mm]
    \foreach \x in {0,...,3}
      \node[note] at (\x, 0.9) {\x};
    \node[note,anchor=east] at (-0.7, 0.9) {value};
    \foreach \v/\a/\b/\c/\d in {0/0/0/0/0, 1/0/0/0/1, 2/1/0/0/1,
        3/1/0/1/1, 4/1/0/1/2, 5/1/1/1/2} {
      \node[note,anchor=east] at (-0.7, -\v) {version \v};
      \node[cell] at (0, -\v) {\a};
      \node[cell] at (1, -\v) {\b};
      \node[cell] at (2, -\v) {\c};
      \node[cell] at (3, -\v) {\d};
    }
    \foreach \v/\i/\x/\c in {1/0/3/1, 2/1/0/1, 3/2/2/1, 4/3/3/2, 5/4/1/1} {
      \node[cell,pb] at (\x, -\v) {\c};
      \node[note,anchor=west] at (3.7, -\v) {adds \(a_{\i} = \x\)};
    }
  \end{scope}
\end{tikzpicture}
\end{center}

Every version is built over the same blocks, so their nodes match up one to one. Subtracting version~\(l\) from version~\(r\) node by node removes the elements before \(l\) from every count. What is left counts exactly the values in \([l, r)\).

\begin{center}
\begin{tikzpicture}[tree,
    level 1/.style={sibling distance=20mm},
    level 2/.style={sibling distance=10mm}]
  \node[tn] at (0, 0) {4}
    child {node {1}
      child {node[label={[note]below:0}] {1}}
      child {node[label={[note]below:1}] {0}}}
    child {node {3}
      child {node[label={[note]below:2}] {1}}
      child {node[label={[note]below:3}] {2}}};
  \node[note] at (0, -34mm) {version 4};
  \node[font=\Large,text=muted] at (25mm, -11mm) {\(-\)};
  \node[tn] at (50mm, 0) {1}
    child {node {0}
      child {node[label={[note]below:0}] {0}}
      child {node[label={[note]below:1}] {0}}}
    child {node {1}
      child {node[label={[note]below:2}] {0}}
      child {node[label={[note]below:3}] {1}}};
  \node[note] at (50mm, -34mm) {version 1};
  \node[font=\Large,text=muted] at (75mm, -11mm) {\(=\)};
  \node[tn] at (100mm, 0) {3}
    child {node {1}
      child {node[label={[note]below:0}] {1}}
      child {node[label={[note]below:1}] {0}}}
    child {node {2}
      child {node[label={[note]below:2}] {1}}
      child {node[label={[note]below:3}] {1}}};
  \node[note] at (100mm, -34mm) {the values in \([1, 4)\)};
\end{tikzpicture}
\end{center}
The small numbers under the leaves are values now, not positions. The difference has one 0, one 2 and one 3, which are exactly \texttt{0 2 3} from \([1, 4)\).

Now walk down the difference from the root. At each node, let \(c\) be the count of its left half. If \(k < c\), the answer is among the smaller values, so go left. Otherwise the left half holds the \(c\) smallest elements, so skip them, subtract \(c\) from \(k\) and go right. The leaf where the walk ends is the answer. Let's find \(k = 1\) in \([1, 4)\).

\begin{center}
\begin{tikzpicture}[tree,
    level 1/.style={sibling distance=20mm},
    level 2/.style={sibling distance=10mm},
    gone/.style={draw=muted,densely dashed,text=muted},
    skip/.style={edge from parent/.style={draw=muted,densely dashed}},
    k/.style={label={[note,text=warncolor]right:\(k = #1\)}}]
  \node[tn,hot,k=1] at (0, 0) {3}
    child {node[pb] {1}
      child {node[label={[note]below:0}] {1}}
      child {node[label={[note]below:1}] {0}}}
    child {node {2}
      child {node[label={[note]below:2}] {1}}
      child {node[label={[note]below:3}] {1}}};
  \node[note] at (0, -34mm) {at the root};
  \node[tn] at (50mm, 0) {3}
    child[skip] {node[gone] {1}
      child {node[gone,label={[note]below:0}] {1}}
      child {node[gone,label={[note]below:1}] {0}}}
    child {node[hot,k=0] {2}
      child {node[pb,label={[note]below:2}] {1}}
      child {node[label={[note]below:3}] {1}}};
  \node[note] at (50mm, -34mm) {at block \([2, 4)\)};
  \node[tn] at (100mm, 0) {3}
    child[skip] {node[gone] {1}
      child {node[gone,label={[note]below:0}] {1}}
      child {node[gone,label={[note]below:1}] {0}}}
    child {node {2}
      child {node[hot,label={[note]below:2}] {1}}
      child {node[label={[note]below:3}] {1}}};
  \node[note] at (100mm, -34mm) {at leaf 2, the answer};
\end{tikzpicture}
\end{center}
At the root \(k = 1\), and the amber left half holds 1 element. Since \(k\) is not below 1, we skip that element and go right with \(k = 0\). At block \([2, 4)\) the left half holds 1 element and \(0 < 1\), so we go left and land on leaf 2. The answer is 2, the middle of \texttt{0 2 3}.

The difference tree is never built. The walk moves down version~\(r\) and version~\(l\) side by side and subtracts their counts at every step. Each step goes one level down, so the whole walk takes \(\bigO{\log m}\).

\subsubsection{The code}
Here is the notebook code piece by piece. The excerpts keep the line numbers of the notebook file. The nodes and the storage come first.

\excerpt[PersistentSegTree]{../notebook/code/ds/persistent-segtree.cpp}{1}{5}

\begin{notes}[Reading the code]
\begin{itemize}
\item Line 2. Children \code|l| and \code|r| are indices into \code|t|, and \code|s| is the sum of the block.
\item Line 3. The array starts with node~0, whose children are node~0 and whose sum is 0. So version~0 is root~0 and nothing is built.
\item Lines 4--5. The tree covers \([0, n)\). When it counts values, \(n\) is the number of values \(m\).
\end{itemize}
\end{notes}

An update copies the path, exactly as drawn.

\excerpt[add]{../notebook/code/ds/persistent-segtree.cpp}{6}{17}

\begin{notes}[Reading the code]
\begin{itemize}
\item Line 7. Start with the whole block \([0, n)\) and return the root of the new version. The old root \code|v| is never touched.
\item Line 9. The copy will sit at the end of \code|t|, at index \code|u|.
\item Line 10. The copy starts out equal to node \code|v|, so it already points at both old children. A vector may copy one of its own elements this way, even when it has to move.
\item Line 11. Every node on the path has \(i\) in its block, so its sum grows by \(x\).
\item Line 12. A leaf has no children to replace.
\item Lines 13--15. Only the child on the side of \(i\) is replaced, by the copy made one level down. The other child stays the old one, and that is the sharing.
\item Line 16. Return the copy, so the parent can point at it.
\end{itemize}
\end{notes}

A query is the plain segment tree query, started at the root of a version.

\excerpt[query]{../notebook/code/ds/persistent-segtree.cpp}{18}{28}

\begin{notes}[Reading the code]
\begin{itemize}
\item Line 23. A block outside \([l, r)\) adds nothing. Node~0 adds nothing either, since its whole block is zeros, so the walk stops there early.
\item Line 24. A block inside \([l, r)\) adds its sum.
\item Lines 25--27. Otherwise split the block and ask both halves.
\end{itemize}
\end{notes}

The \(k\)th smallest walks down two versions side by side.

\excerpt[kth]{../notebook/code/ds/persistent-segtree.cpp}{29}{39}

\begin{notes}[Reading the code]
\begin{itemize}
\item Line 30. Pass \code|a| as the root of version~\(l\) and \code|b| as the root of version~\(r\), so the walk counts the elements of \([l, r)\).
\item Line 31. \code|[lo, hi)| is the block of values under the current pair of nodes.
\item Line 34. \code|c| is the count of the left half in the difference, the left child of \code|b| minus the left child of \code|a|.
\item Line 35. The answer is among the smaller values, so both nodes go left.
\item Line 36. Skip the \code|c| smaller elements, and both nodes go right.
\item Line 38. A single value is left, and it is the answer. It is a rank, so map it back to the real value.
\end{itemize}
\end{notes}

In a problem, the values are replaced by their ranks first, and the roots go in a list.

\begin{cpp}[listing options={style=cp,emph={PersistentSegTree,add,kth}}]
vi vals = a;  // the distinct values, sorted
sort(all(vals));
vals.erase(unique(all(vals)), vals.end());
PersistentSegTree st(sz(vals));
vi root{0};  // root[i] counts a[0, i)
for (int x : a) {
  int y = lower_bound(all(vals), x) - vals.begin();
  root.push_back(st.add(root.back(), y, 1));
}
// k-th smallest of a[l, r), k from 0
int ans = vals[st.kth(root[l], root[r], k)];
\end{cpp}

\begin{pitfall}
Never keep a reference into \code|t| across a call to \code|add|, because \code|add| grows the vector and it may move. Lines 14 and 15 are safe only because C++17 evaluates the right side of an assignment first. Each update adds up to \(\ceil{\log_2 n} + 1\) nodes of 16 bytes, so \(2 \cdot 10^5\) updates take about 60~MB. And \code|kth| needs \(k < r - l\), or it ends on a wrong leaf.
\end{pitfall}
